\providecommand{\Sp}{\mathrm{Sp}}
\providecommand{\Mp}{\mathrm{Mp}}
\providecommand{\Z}{\mathbb{Z}}
\providecommand{\Q}{\mathbb{Q}}
\providecommand{\C}{\mathbb{C}}
\providecommand{\R}{\mathbb{R}}
\providecommand{\CY}{\mathrm{CY}}
\providecommand{\SVOA}{\mathrm{SVOA}}
\providecommand{\VOA}{\mathrm{VOA}}
\providecommand{\NS}{\mathrm{NS}}
\providecommand{\ChirAlg}{\mathrm{ChirAlg}}
\providecommand{\Cat}{\mathrm{Cat}}
\providecommand{\cC}{\mathcal{C}}
\providecommand{\Coh}{\mathrm{Coh}}
\providecommand{\cO}{\mathcal{O}}
\providecommand{\Aut}{\mathrm{Aut}}
\providecommand{\id}{\mathrm{id}}
\providecommand{\kBKM}{\kappa_{\mathrm{BKM}}}
\providecommand{\fg}{\mathfrak{g}}
\providecommand{\fsg}{\mathfrak{sg}}
\providecommand{\osp}{\mathfrak{osp}}
\providecommand{\Phisup}{\Phi^{\mathrm{sup}}}

\section{Metaplectic CY\(_2\) data and the super vertex functor}
\label{sec:phi-sup-functor}

\subsection*{The source category \(\CY_2^{\Mp}\)}

The source category \(\CY_2^{\Mp}\) consists of pairs
\((\cC,\nu)\) with \(\cC\in\CY_2\text{-}\Cat\) smooth proper and
\[
  \nu:\pi_1(\mathrm{Per}(\cC))\longrightarrow \Mp_4(\Z)
\]
a metaplectic lift of the period-map monodromy. Equivalently, \(\nu\)
is a theta multiplier on the period stack cover, compatible with the
Mukai pairing on \(H^*(\cC;\Z)_{\mathrm{alg}}\). Morphisms are
CY-Fourier-Mukai kernels whose induced Mukai isometries intertwine the
chosen multipliers up to Weil-cocycle coboundary (Shimura 1975
\emph{Ann.\ Math.}\ \textbf{102}; Howe 1979
\emph{Proc.\ Symp.\ Pure Math.}\ \textbf{33}). The forgetful map
\(\CY_2^{\Mp}\to\CY_2\) is a \(\Z/2\)-gerbe, and the trivial
multiplier locus embeds by \(\cC\mapsto(\cC,\id)\).

\subsection*{The target category \(\SVOA\)}

Objects of \(\SVOA\) are \(\Z/2\)-graded vertex superalgebras
\[
  V=V_{\bar 0}\oplus V_{\bar 1}
\]
equipped with a Neveu-Schwarz super-Virasoro subalgebra
\(\NS_{c^{\mathrm{sup}}}\subset V\). In the FMS normalisation,
\[
  c^{\mathrm{sup}}(V)
  =
  c_{\bar 0}(V)+\frac12 c_{\bar 1}(V),
\]
where \(c_{\bar 0}\) is the even stress-tensor contribution and
\(c_{\bar 1}\) counts odd \(\beta\gamma\)-\(bc\) first-order pairs
after Friedan-Martinec-Shenker bosonisation (FMS 1986
\emph{Nucl.\ Phys.\ B}~\textbf{271}). Morphisms are parity-preserving
vertex-superalgebra maps. The full subcategory with \(V_{\bar 1}=0\)
is the ordinary VOA subcategory (Kac 1998, \S5.9).

\subsection*{The assignment \(\Phisup\)}

The functor is defined on the decorated subcategory
\(\widetilde{\CY}_2^{\Mp}\) whose objects are \((\cC,\nu)\) together
with the data
\[
\begin{gathered}
  \Lambda_{(\cC,\nu)}\subset H^*(\cC;\Z)_{\mathrm{alg}}
  \quad\text{even, integral, and \(\nu\)-invariant},\\
  \fsg_{(\cC,\nu)}
  \quad\text{a BKM superalgebra with denominator \(\Phi_\nu\)},\\
  r_f(\nu)
  \quad\text{the number of odd imaginary simple-root generators}.
\end{gathered}
\]
The denominator \(\Phi_\nu\) is the Gritsenko-Borcherds lift of the
chosen twined elliptic genus, and the root multiplicities are read from
its Fourier coefficients. The metaplectic lift supplies the multiplier;
the lattice, parity, and BKM presentation are additional domain data.
Set
\[
  \Phisup(\cC,\nu)
  =
  V_{\Lambda_{(\cC,\nu)}}^{\bar 0}
  \otimes
  V_{\mathrm{FMS}}^{\bar 1}\bigl(r_f(\nu)\bigr).
\]
The first factor is the lattice VOA attached to the even lattice
\(\Lambda_{(\cC,\nu)}\) (Borcherds 1986 \emph{PNAS}~\textbf{83}); the
second factor is the FMS vertex-superalgebra for the odd root sector.

On a Fourier-Mukai kernel
\[
  f:(\cC,\nu)\longrightarrow(\cC',\nu')
\]
that preserves the parity datum and transports the Weil cocycle, the
bosonic component of \(\Phisup(f)\) is the lattice-VOA map induced by
the Mukai isometry
\(\Lambda_{(\cC,\nu)}\to\Lambda_{(\cC',\nu')}\). The fermionic
component is the FMS action on the corresponding \(\Z/2\)-graded
Dirac modules. Thus \(\Phisup\) is a functor on the subcategory whose
morphisms preserve all three displayed data.

\subsection*{Restriction to the trivial multiplier}

On the \(\Mp\)-trivial section where the \(\Phi_2\)-output is realised
by the same lattice VOA construction, the odd FMS factor vanishes and
\[
  \Phisup(\cC,\id)=\Phi_2(\cC)
\]
as an \(\SVOA\)-object concentrated in even parity.
The comparison square is then
\[
\begin{array}{ccc}
  \CY_2 & \xrightarrow{\Phi_2} & E_2\text{-}\ChirAlg \\
  \downarrow & & \downarrow \\
  \CY_2^{\Mp} & \xrightarrow{\Phisup} & \SVOA .
\end{array}
\]
This is a CY\(_2\) statement. The specialised CY\(_3\) output remains
\(E_1\).

\subsection*{The \(K3\times E\) invariants}

For \(X=K3\times E\), Kunneth gives
\[
  H^{0,q}(X)
  \cong
  \bigoplus_{a+b=q}H^{0,a}(K3)\otimes H^{0,b}(E),
\]
with \(h^{0,0}(K3)=h^{0,2}(K3)=1\), \(h^{0,1}(K3)=0\), and
\(h^{0,0}(E)=h^{0,1}(E)=1\). Thus the compact chiral Hodge supertrace is
\[
  \kappa_{\mathrm{ch}}(A_X)
  =
  \sum_{q=0}^{3}(-1)^q h^{0,q}(X)
  =
  1-1+1-1
  =
  0.
\]
The categorical Euler characteristic is Kunneth-multiplicative:
\[
  \kappa_{\mathrm{cat}}(X)
  =
  \chi(\cO_{K3})\chi(\cO_E)
  =
  2\cdot0
  =
  0.
\]
The Heisenberg-Mukai specialisation is a different chiral
algebraisation:
\[
  \kappa_{\mathrm{ch}}^{\mathrm{Heis}}(X)=2+1=3.
\]
The Borcherds branch reads the weight of the level-one denominator:
\[
  \kBKM(\Delta_5)=\frac{c_1(0)}2=\frac{10}{2}=5.
\]
The fibre invariant is the Mukai-lattice rank
\[
  \kappa_{\mathrm{fiber}}(X)=\operatorname{rk}H^*(K3;\Z)=24.
\]
These values are attached to different constructions. The compact
Hodge branch, the Heisenberg branch, the Borcherds branch, and the
Mukai-fibre branch are kept separate throughout the comparison.

\subsection*{Compatibility with compact CHL CY\(_3\) data}

For a compact CHL orbifold
\[
  X^{(N)}=(K3\times E)/\Z_N,\qquad N\in\{1,2,3,4,6\},
\]
the specialised CY\(_3\) output of the two-stage CY-to-chiral
construction lies in \(E_1\text{-}\ChirAlg\):
\[
\begin{array}{ccc}
  \CY_3 & \xrightarrow{\Phi_3} & E_1\text{-}\ChirAlg \\
  \uparrow\pi_* & & \uparrow\mathrm{bdy} \\
  \CY_2^{\Mp} & \xrightarrow{\Phisup} & \SVOA .
\end{array}
\]
A bulk-to-boundary comparison on the CHL locus requires a near-horizon
K3 slice with metaplectic multiplier \(\nu_N\), a boundary VOA
construction for the holomorphic twist, and a compatibility map from
the \(E_1\)-chiral object to the super-Virasoro module. Costello-Gaiotto's
boundary VOA construction supplies the model for this map in
the holomorphic-twist setting. An \(E_2\)-object appears only after a
center, Drinfeld-double, module-category, or boundary-enhancement
construction with its Hall hypotheses.

The BKM branch in this paragraph is the Hall-Drinfeld/Borcherds branch
for \(K3\times E\). The Mukai self-mirror K3 Yangian branch
\(Y(\mathfrak g_{K3})\) on \(\Lambda_{\mathrm{Muk}}\) is a separate
CY\(_2\) integrable construction.

\subsection*{Borcherds constants}

The primitive compact CHL levels are indexed by
\[
  N\in\{1,2,3,4,6\}.
\]
For the primitive denominator \(\Phi_N\), the zero-mode constant and
the Borcherds weight are
\[
\begin{array}{c|ccccc}
N&1&2&3&4&6\\
\hline
c_N(0)&10&8&6&4&2\\
\kBKM(\Phi_N)=c_N(0)/2&5&4&3&2&1
\end{array}
\]
in the Gritsenko-Nikulin CHL normalisation.

The Gritsenko-Cl\'ery diagonal-divisor catalogue is a separate
eight-row paramodular catalogue. Its constants are
\[
\begin{array}{c|cccccccc}
j&1&2&3&4&5&6&7&8\\
\hline
c_j^{\mathrm{dd}}(0)&10&4&6&2&4&1&3&2\\
\kBKM(F_j)=c_j^{\mathrm{dd}}(0)/2&5&2&3&1&2&1/2&3/2&1
\end{array}
\]
The two tables share only the level-one row \(\Delta_5\) without an
additional row-map datum.

Rows \(j=6,7,8\) of the Gritsenko-Cl\'ery catalogue are metaplectic
period-side rows:
\[
\begin{array}{c|c|c|c}
j & c_j(0) & \kBKM(\Phi_j)=c_j(0)/2 & \text{cover/multiplier} \\
\hline
6 & 1 & 1/2 & M_{1/2}(\Gamma_4,\nu_8) \\
7 & 3 & 3/2 & M_{3/2}(\Gamma_1(4),\nu_4) \\
8 & 2 & 1 & M_{1}(\Gamma_6,\nu_6)
\end{array}
\]
These constants determine the Borcherds weights. The super-VOA output
requires the additional lattice, parity, and FMS data specified above.
For such a datum,
\[
  \Phisup(K3,\nu_j)
  =
  V_{\Lambda_j}^{\bar 0}
  \otimes
  V_{\mathrm{FMS}}^{\bar 1}(r_f(j)),
  \qquad
  c^{\mathrm{sup}}_j=\mathrm{rk}(\Lambda_j)+\frac12 r_f(j).
\]

\subsection*{Metaplectic super-vertex sector}

Compact CHL levels give the integer \(\kBKM\) locus coming from compact
smooth CY\(_3\) quotients. The metaplectic rows enlarge the period-side
CY\(_2\) input by allowing half-integral Borcherds weights on
\(\Mp_4\)-covers. The extraction rule
\[
  \kBKM(\Phi)=c(0)/2
\]
governs both indexing families after the relevant row datum is fixed.
Compact CY\(_3\) specialisation gives \(E_1\)-chiral output. The
metaplectic CY\(_2\) construction gives a super-VOA after the
lattice, FMS, and Sugawara data are fixed.

A \(Y_{\osp(4|20)}\) Super-Yangian requires generators, relations,
coproduct, integral form, reflection-equation data, and comparison with
the Hall-Drinfeld/Borcherds branch. The construction above supplies the
decorated super-VOA object; the Super-Yangian remains conjectural.

\paragraph{Primary literature.}
Kac 1998, \emph{Vertex Algebras for Beginners};
Kac-Wakimoto super-affine (Kac 1977 \emph{Adv.\ Math.}~\textbf{26};
Kac-Wakimoto 1989 \emph{Proc.\ Nat.\ Acad.\ Sci.}~\textbf{85});
Friedan-Martinec-Shenker 1986 \emph{Nucl.\ Phys.\ B}~\textbf{271};
Borcherds 1986 \emph{PNAS}~\textbf{83};
Shimura 1975 \emph{Ann.\ Math.}~\textbf{102};
Howe 1979 \emph{Proc.\ Symp.\ Pure Math.}~\textbf{33};
Gritsenko-Nikulin 1998 \emph{Amer.\ J.\ Math.}~\textbf{119};
Cl\'ery-Gritsenko 2008 \emph{Acta Arith.}~\textbf{133};
Cheng-Duncan-Harvey 2014 \emph{Comm.\ Num.\ Th.\ Phys.}~\textbf{8};
Costello-Gaiotto 2018;
Lorgat 2020 (arXiv:2004.01676).
